\chapter{Design Pattern}

Cette partie du projet vise à améliorer la qualité et la réutilisabilité du code ITIAventure en appliquant le Design Pattern Abstract Factory. Ce patron permet de centraliser la création des entités du jeu dans des classes dédiées, tout en facilitant l’extension à d’autres univers sans modifier le code existant.

\section{ITIAventureFactory}

ITIAventureFactory regroupe les méthodes de création des entités de base : Monde, Piece, Porte, Monstre, Serrure, Cle.
Elle encapsule les appels aux constructeurs, ce qui permet de créer des entités sans connaître leur classe concrète.

\section{ITISpaceOperaFactory}

Pour le nouvel univers spatial, nous avons créé ITISpaceOperaFactory, qui hérite de ITIAventureFactory et redéfinit les méthodes de création pour retourner des entités spécifiques : Galaxie, VaisseauSpatial, Teleporteur, Alien, LecteurBadge, Badge.

Grâce à ce système :

    le code principal du jeu reste inchangé ;

    on peut facilement basculer entre les univers ;

    les classes restent ouvertes à l’extension mais fermées à la modification.

\section{Conclusion}

L’usage du pattern Abstract Factory nous a permis de découpler la logique de création des entités du reste du jeu.